\section{Introduction}\label{sec:intro}

The sums of at most $h$ elements of a finite set of integers need not fill the range below $h$ times its largest element. This paper measures how far they fall short: for each target, the distance down to the nearest sum not above it, added over all targets, as a function of $h$. Let $A_0=\{0=a_0<a_1<\dots<a_k=d\}$ be a set of integers with $\gcd(A_0)=1$, and let its $h$-fold sumset $hA_0\subseteq[0,hd]$ be the set of sums of at most $h$ elements of $A=A_0\setminus\{0\}$. For a target $T\ge0$ the \emph{undershoot}
\[
D_h(T)=T-\max\{n\in hA_0:n\le T\}
\]
is the distance from $T$ down to the largest sum not above it, and $\Tot(h)=\sum_{T=0}^{hd}D_h(T)$ is its total over the targets in range. The question came from filling advertising breaks: if the elements of $A$ are the available ad durations and a break of length $T$ takes at most $h$ ads, then $D_h(T)$ is the time left unfilled when the break is filled as fully as possible. The range of targets grows with $h$, yet the total stays bounded and is eventually constant. How it reaches its final value depends on the set. For $\{0,5,7,14\}$ the totals at $h=1,2,\dots$ are $32,48,56,61,64,65,65,\dots$: they rise by smaller steps each time and then stop. For $\{0,2,8,11\}$ they are $19,19,18,18,\dots$: they start above the final value and come down to it. For $\{0,4,9,13\}$ they are $22,36,45,48,49,46,44,40,40,\dots$: they climb past the final value $40$ and return to it only at $h=8$. This paper explains the three behaviors and describes the whole sequence $\Tot(h)$.

\paragraph{The limit.} For large $h$ the set $hA_0$ is an interval with fixed sets of exceptions at both ends \cite{Nathanson1972,GranvilleShakan2020,GranvilleWalker2021,Lev2022}, governed by the two numerical semigroups $S_1=\langle A\rangle$ and $S_2=\langle d-A_0\rangle$. Through the projective monomial curve of $A_0$, the level from which this holds is a regularity \cite{Elias2022,GimenezGonzalezSanchez2023}, and the sizes $|hA_0|$ have recently been computed, for special classes of sets \cite{Nathanson2025} and, from the strict records of the least number of parts in each residue class, for every finite set \cite[Thm.~1]{ChenYang2026}. These results count the sums. The undershoot records how far each target is from one. The end structure gives its limit: for every $A_0$ and all large $h$, $\Tot(h)=\M(S_1)+\M(S_2)$ (Propositions~\ref{prop:failing-interval} and~\ref{prop:overshoot}), where the mass $\M(S)$ of a numerical semigroup $S$ is the sum of $L(L+1)/2$ over its deserts, the maximal runs of gaps, of lengths $L$. For the three sets above the limits are $65$, $18$ and $40$. The end structure does not show how the totals approach the limit, and that is the subject of this paper.

\paragraph{Main results.} Let $c_1$ and $c_2$ be the conductors of $S_1$ and $S_2$, and let
\[
I_h=\{n\in[0,hd]:n\in S_1,\ hd-n\in S_2\}.
\]
Always $hA_0\subseteq I_h$. We call $A_0$ \emph{Cohen--Macaulay} (CM) when $hA_0=I_h$ for every $h$; by Cavaliere and Niesi this is the arithmetic Cohen--Macaulay property of the monomial curve (Theorem~\ref{thm:cm}). For a CM set and each residue $s$ modulo $d$ there are integers $u_s,v_s\ge0$, read off the Ap\'ery sets of $S_1$ and $S_2$, such that $kd-s\in hA_0$ exactly when $u_s\le k\le h-v_s$ (Lemma~\ref{lem:residue}). We call the pairs $(u_s,v_s)$ \emph{similarly ordered} if $u_s>u_t$ implies $v_s\ge v_t$, the hypothesis of Chebyshev's sum inequality. The sumsets regularity $\sigma(A)$ of Gimenez and Gonz\'alez-S\'anchez is the least level from which $hA_0$ keeps its final shape (Section~\ref{sec:prelim}).

\begin{maintheorem}[Theorem~\ref{thm:spectrum}, Proposition~\ref{prop:threshold}]\label{thm:A}
Let $A_0$ be CM with $d\ge2$. There are integers $\lambda_1\ge\lambda_2\ge\dots\ge0$, finitely many of them positive, with $\sum_\ell\lambda_\ell=\M(S_1)+\M(S_2)$, such that
\[
\Tot(h)\ge\lambda_1+\dots+\lambda_h\qquad(h\ge1),
\]
with equality for every $h$ if and only if the pairs $(u_s,v_s)$ are similarly ordered. In that case $\Tot$ is nondecreasing and concave, $\Tot(h)=\M(S_1)+\M(S_2)$ exactly when $h\ge\sigma(A)$, and the Castelnuovo--Mumford regularity of the coordinate ring of the curve is $\sigma(A)+1$.
\end{maintheorem}

Proposition~\ref{prop:threshold} proves the threshold and the regularity statement at every modulus $m\ge2$, a divisor of $d$ lying in $S_1\cap S_2$, at which the pairs read modulo $m$ are similarly ordered. The set $\{0,5,7,14\}$ is similarly ordered with $\lambda=(32,16,8,5,3,1)$, and its totals are the partial sums. The set $\{0,2,8,11\}$ is not: its bound $10,15,18$ is strict at $h=1$ and $h=2$, and the levels at which the bound is strict are given in general by the inversions of the pairs (Proposition~\ref{prop:strict-levels}). For four elements similar ordering has a finite geometric test. The pairs of a CM set $\{0,a_1,a_2,d\}$ come from its Ap\'ery diagram, a set of $d$ points $(x,y)\in\N^2$ such that the values $a_1x+a_2y$ are the least elements of $S_1$ in the $d$ residue classes (Lemma~\ref{lem:diagram}).

\begin{maintheorem}[Theorem~\ref{thm:similar}, Proposition~\ref{prop:shape}]\label{thm:B}
For a CM set $\{0,a_1,a_2,d\}$ the Ap\'ery diagram is a rectangle or an L-shape, and the pairs are similarly ordered if and only if on each layer $x+y=t$ of the diagram the values $a_1x+a_2y$ lie in one interval $((q-1)d,qd]$ with $q\in\Z$. An L-shape whose notch has both sides of length at least $2$ is never similarly ordered.
\end{maintheorem}

The family $\{0,a,b,2b\}$ is CM. Its pairs need not be similarly ordered at $d=2b$, but they always are when the residues are taken modulo $b$, and its spectrum is then a statistic of an arithmetic progression on a circle. For $1\le j\le b$ let $\phi(j)$ be the sum of $\binom g2$ over the gaps $g$ of the points $0,a,\dots,(j-1)a$ on $\Z/b\Z$.

\begin{maintheorem}[Theorems~\ref{thm:family-spectrum} and~\ref{thm:three-distance}, Corollaries~\ref{cor:family-mass} and~\ref{cor:quadratic-pieces}]\label{thm:C}
Let $0<a<b$ be coprime and $A_0=\{0,a,b,2b\}$. Then
\[
\Tot(h)=\binom b2+2\sum_{j=2}^{h}\phi(j)+\phi(h+1)\qquad(1\le h\le b-1),
\]
and $\Tot(h)=2\sum_{j=1}^b\phi(j)-\binom b2$ for $h\ge b-1$. The gaps of $\{0,a,\dots,(j-1)a\}$ take at most three lengths, given by the continued fraction of $a/b$ through the three-distance theorem, and $\Tot$ consists of explicit quadratic pieces, one for each block of that continued fraction.
\end{maintheorem}

For $\{0,5,7,14\}$, $\phi(1),\dots,\phi(6)=21,11,5,3,2,1$, which gives $32,48,56,61,64,65$. A set that is not CM has \emph{failing levels}, the $h$ with $hA_0\neq I_h$. They are finitely many, and $\tau$ denotes the level after the last one.

\begin{maintheorem}[Propositions~\ref{prop:failing-interval} and~\ref{prop:overshoot}, Corollary~\ref{cor:four-overshoot}, Example~\ref{ex:five-element}]\label{thm:D}
Let $A_0$ be arbitrary and $hd\ge c_1+c_2$. Then $\Tot(h)\ge\M(S_1)+\M(S_2)$, with equality if and only if $h$ is not a failing level. Every four-element set that is not CM has $\Tot(\tau-1)>\M(S_1)+\M(S_2)$. The five-element set $\{0,5,9,11,20\}$ is not CM, and $\Tot(h)\le\M(S_1)+\M(S_2)$ for every $h$.
\end{maintheorem}

The set $\{0,4,9,13\}$ has $c_1+c_2=48$ and failing levels $4$ to $7$, so by Theorem~\ref{thm:D} its total exceeds $40$ at each of them and equals $40$ from $h=8$ on. The four-element statement uses the companion paper \cite[Theorem~A]{CompanionRegularity}, which shows that $(\tau-1)d\ge c_1+c_2+d$ for these sets. For arbitrary sets the undershoot differs from that of $I_h$ by an explicit correction (Proposition~\ref{prop:correction}), and finitely many corners per residue class give the exact undershoot over any window of targets (Theorem~\ref{thm:hinge}).

\paragraph{Method.} Everything is read off the residue pairs $(u_s,v_s)$, which locate each residue class in $hA_0$ block by block. For a CM set the integers of $[0,hd]$ missing from $hA_0$ fall in blocks between consecutive multiples of a modulus, and the undershoot of a target is the number of consecutive missing integers ending at it (Proposition~\ref{prop:counting}). The missing positions of a block are the union of a set coming from $S_1$ and a set coming from $S_2$. Taking the larger of the two in every block gives the spectral bound, and the bound is attained in a block exactly when the two sets are nested. At the modulus $d$ they are nested in every block at every level exactly under similar ordering. For four elements the pairs come from the Ap\'ery diagram, and the test becomes a test on its layers. For $\{0,a,b,2b\}$ the missing positions at the modulus $b$ are the complements of the progressions $0,-a,\dots,-(j-1)a$ on $\Z/b\Z$, so the spectrum is a sum over circular gaps and the three-distance theorem evaluates it. For a set that is not CM we compare $hA_0$ with $I_h$: once $hd\ge c_1+c_2$ the total of $I_h$ is the mass, and every element of $I_h$ missing from $hA_0$ raises the total. The corrections and window formulas come from the minimal corners of the levels at which each class is present.

\paragraph{Related work.} The three-distance theorem is due to S\'os \cite{Sos1958}, Sur\'anyi \cite{Suranyi1958} and \'Swierczkowski \cite{Swierczkowski1958}; we use it in the explicit form of Alessandri and Berth\'e \cite{AlessandriBerthe1998}. The homogenized Ap\'ery set and the Cohen--Macaulay criterion are those of Cavaliere and Niesi \cite{CavaliereNiesi1983} and Gimenez and Gonz\'alez-S\'anchez \cite{GimenezGonzalezSanchez2023}. Lemma~\ref{lem:residue} is their coordinate description read on one level, and Proposition~\ref{prop:counting} follows from it directly. Proposition~\ref{prop:failing-interval} is \cite[Lemma~2.5]{CompanionRegularity} in other coordinates, and the L-shapes of Section~\ref{sec:similar} are those of three-generated numerical semigroups \cite{AguiloMarijuan2014}. What this paper adds begins with the spectrum. Hoa and Tien \cite{HoaTien2026} bound the regularity of the monomial curve through Ap\'ery sets and Frobenius numbers, with criteria for the Cohen--Macaulay property and an application to the structure of sumsets. Similar ordering is the condition of Chebyshev's sum inequality \cite{HardyLittlewoodPolya1952}. To our knowledge the undershoot has not been studied as a functional of $hA_0$. The main question left open is a description of the similarly ordered four-element sets by a short list of congruence or continued-fraction conditions on $(a_1,a_2,d)$ (Section~\ref{sec:open}).

\paragraph{Organization.} Section~\ref{sec:prelim} fixes notation. Sections~\ref{sec:sumsets} to~\ref{sec:noncm} prove Theorems~\ref{thm:A} to~\ref{thm:D} in order, and Section~\ref{sec:corrections} gives the corrections and window formulas. Section~\ref{sec:open} lists open problems, and Section~\ref{sec:computation} describes the computations, all exact and reproducible from the accompanying code \cite{code2026}.
